% Input by mostlyharmless.tex.  No preamble here.

\section{The two interactive models}
\label{sec:models}

The derivations are complete without this section, and nothing here is certified: what
follows is a pair of front ends and the visual checks that go with them.  They are recorded
because a picture can be wrong in ways a residual cannot --- a mirrored image satisfies every
statement in sections~\ref{sec:visibility} to~\ref{sec:full-map}, and did so for long enough
to be embarrassing --- and because a check that lives in prose is one that drifts away from
the code it describes.

Both models are Matplotlib windows over the same scaffolding: two panels, a shared colour
bar, sliders, radio groups, and a redraw that fires once when a slider drag ends rather than
once per mouse motion.  Each carries a \emph{preset} group: named parameter sets, labelled by
what you should see rather than by their angles, with a \code{manual} entry meaning ``these
are the user's own numbers'', re-selected whenever a slider moves.  The presets are the
reproducible part of the checks below.

\subsection{The deformation model}
\label{sec:model-deformation}

\code{make gui} shows a rotationally deformed star in the far field: a unit-radius source
whose level surface is the equipotential of a tide and a spin, with brightness fixed at the
von Zeipel law $I \propto g$.  The two panels are the near and far hemispheres.  The sliders
are the two amplitudes, the inclination, and the tide's body latitude and longitude; the
position angle of the spin is pinned to zero, being a gauge for a single observation.

The row group is the part worth explaining, because the rows differ in kind:
\begin{description}
\item[\code{numeric}] the exact level surface, solved by bisection along each sampled
  direction.  The reference.
\item[\code{first-order}] an ellipsoid, with the brightness linearised as in
  \eqref{eq:first-order-pattern}.  Against \code{numeric} this measures the whole analytic
  stack at $O(\epsilon)$.
\item[\code{fitted, l = 0,2,4}] coefficients chosen by least squares in the image plane,
  rather than taken from the brightness pattern.
\item[\code{fitted, l = 0,2,4,6}] the same at one more degree.  Comparing the two fitted rows
  isolates the \emph{truncation} from the shape model's error, which the other pairs
  conflate.
\end{description}
The fitted rows print their residual in the panel titles.  It is not a defect of the fit but
the price of the degree: at fixed $\Lambda$ the residual is the truncation floor, and
section~\ref{sec:full-map} states why a discontinuous map converges slowly.

\paragraph{The presets, and what they test.}  \code{reference} is a moderate configuration,
$\epsilon_\text{tide} = 0.03$ and $\epsilon_\text{spin} = 0.01$, keeping
$\epsilon_\text{tide} + \tfrac32\epsilon_\text{spin} \approx 0.045$ well below the $0.148$
ceiling, with the spin axis inclined by $60°$ so that both the tide's bulge and the spin's
flattening are visible.  \code{ceiling} puts the same amplitudes in the equatorial plane,
which is the stress test for the analytic rows: the affine error reaches $0.11$ there.

\code{tilt check} is the geometry test.  The deformations are at the ceiling and the tide axis
is off both axes of the sky, at latitude $-45°$ and longitude $45°$, so the bulge lies on a
\emph{diagonal}.  That matters: an axis-aligned bulge maps to itself under a mirror, so a
sign error in the frame or in the display would pass unnoticed, while a diagonal one would
come out on the wrong diagonal.  Seen on the page, the stretching should run from lower right
to upper left --- mostly North with a little East, with the body-latitude bands
perpendicular to it --- and the picture is drawn with North up and East to the left, as
section~\ref{sec:conventions} fixes.

\subsection{The map model}
\label{sec:model-map}

\code{make mapgui} shows a starspot as spherical-harmonic coefficients: a polar cap carried in
the body frame and rotated into the observer frame, in the manner of
section~\ref{sec:rotation-law}.  There is no deformation and no solve --- the model is a
\code{Coeffs} object --- so a redraw costs a coefficient sum, and the panels are the near and
far hemispheres in units of the mean.

The sliders are the spot's half-angle, its depth, its colatitude away from the spin axis, its
longitude, the inclination of the spin to the line of sight, and the position angle of the
sky.  The colatitude is what makes the longitude mean something: the cap is built about the
pole, and an axisymmetric map is invariant under a rotation about its own axis, so with the
spot at the pole the longitude slider does nothing at all.  The \code{l_max} radio is a
truncation knob, and switching it shows the ringing of
certificate~\vtag{28} directly.

\paragraph{The presets, and what they test.}  Four scenes, each labelled by where the spot
should appear.  With the spin pointing at the observer --- inclination $0$ --- and the spot at
colatitude $60°$:
\begin{description}
\item[\code{spot left}] longitude $90°$: the spot must appear \emph{left} of centre,
  vertically centred.
\item[\code{spot right}] longitude $270°$: the mirror position, \emph{right} of centre.
\item[\code{spot above}] longitude $0$: \emph{above} centre.
\item[\code{spot behind}] inclination $90°$, longitude $0$: the spin axis now lies in the
  plane of the sky, so the spot goes to the \emph{far} panel, and the near panel shows
  nothing of it.
\end{description}
The first three together pin the page: East leftwards and North upwards cannot both be wrong
and still put the spot where the label says.  The fourth exercises the far hemisphere and its
mirror, which no residual can check.  The visibility view adds a fifth, arithmetic check that
needs no preset --- for a real map $V(-u,-v) = \overline{V(u,v)}$, so the real part of
$\Delta V$ is even under $(u,v) \to (-u,-v)$ and the imaginary part odd
\eqref{eq:hermitian}.

\subsection{What the checks are worth}
\label{sec:models-worth}

The two models are deliberately outside the certificate scheme.  A certificate compares two
routes to the same number; the checks here compare a picture with a stated expectation, and
their evidence is that a wrong answer is visible rather than that a residual is small.  That
is weaker, and it is why the presets exist: they move the expectation out of prose and into
the program, so that the geometry can be re-checked without anyone having to remember which
slider to move, or in which direction the answer should come out.
